\documentclass[11pt]{article}

\usepackage{amsmath,amssymb,amsthm}
\usepackage{graphicx}
\usepackage{listings}
\usepackage{url}

\title{AICE: A Typed Actor-based Intelligent Computing Environment}

\author{
Yasushi Kodama\\
Faculty of Business Administration\\
Hosei University
}

\date{}

\begin{document}

\maketitle

\begin{abstract}
We present AICE (Actor-based Intelligent Computing Environment), a lightweight and distributed AI execution framework built on a typed actor language AIPL. AICE integrates asynchronous message passing, selective receive, and AI agent composition under a statically typed setting.

We introduce a three-layer type system consisting of value types, actor types, and communication types, and show that type safety alone is insufficient to guarantee correctness of asynchronous communication.

We formalize the operational semantics of selective receive and propose an optional integration of session types for protocol-level guarantees.
\end{abstract}

\section{Introduction}

Modern AI systems require distributed, concurrent, and lightweight execution environments. Existing systems lack formal guarantees and are often too heavy for edge or embedded deployment.

This paper proposes AICE, a typed actor-based computing environment built on AIPL.

\paragraph{Contributions}
\begin{itemize}
\item A typed actor language AIPL
\item AICE architecture for AI execution
\item Formal operational semantics for asynchronous messaging
\item Analysis of limitations of type safety
\item Optional session type integration
\end{itemize}

\section{Language Overview}

\subsection{Syntax}

\begin{align*}
e &::= x \mid n \mid e_1 + e_2 \mid new\ C \mid send\ e.m(e^*) \\
s &::= e \mid s_1 ; s_2 \mid if\ e\ then\ s_1\ else\ s_2 \mid select\ \{b_i\} \\
b &::= when\ m(x^*) \rightarrow s
\end{align*}

\section{Type System}

\subsection{Types}

\[
\tau ::= int \mid string \mid unit \mid \tau_1 \to \tau_2 \mid actor(C)
\]

Extended form:

\[
actor(C, P)
\]

\subsection{Typing Judgement}

\[
\Gamma \vdash e : \tau
\quad
\Gamma \vdash s : unit
\]

\subsection{Send Typing}

\[
\frac{
\Gamma \vdash e : actor(C)
\quad
m : \tau_1 \times \dots \times \tau_n \to unit \in C
\quad
\Gamma \vdash e_i : \tau_i
}{
\Gamma \vdash send\ e.m(e_1,\dots,e_n) : unit
}
\]

\subsection{Select Typing}

\[
\frac{
\Gamma \vdash s_i : unit \quad \forall i
}{
\Gamma \vdash select\ \{b_i\} : unit
}
\]

\section{Operational Semantics}

\subsection{Configuration}

\[
\langle A, M, s \rangle
\]

\subsection{Send}

\[
\langle A, M, send\ B.m(v) \rangle \rightarrow \langle A, M, unit \rangle
\]

Mailbox update:

\[
MB_B := MB_B \cup \{(m, v, A)\}
\]

\subsection{Select}

\[
\frac{
(m, v, sender) \in M
}{
\langle A, M, select \rangle \rightarrow \langle A, M \setminus \{m\}, s \rangle
}
\]

\subsection{Blocking}

If no message matches:

\[
\langle A, M, select \rangle \rightarrow wait
\]

\subsection{Non-determinism}

Multiple matches:

\[
select \rightarrow s_1 \quad \text{or} \quad s_2
\]

\section{Limitations of Type Safety}

We show that type safety does not guarantee communication correctness.

\subsection{Example}

\begin{lstlisting}
send calc.add(3,4);
send calc.result(10);
\end{lstlisting}

Even if well-typed, execution order is not guaranteed.

\section{Session Types (Optional)}

\subsection{Protocol}

\[
P = A \rightarrow B : m_1 ; B \rightarrow A : m_2
\]

\subsection{Design}

Session types are optional and do not interfere with HM inference.

\section{AICE Architecture}

\begin{itemize}
\item Actor runtime
\item Scheduler
\item Mailbox system
\item Network layer
\item AI actors
\end{itemize}

\subsection{AI Actor Roles}

\begin{itemize}
\item PlannerActor
\item LLMActor
\item MemoryActor
\item ToolActor
\end{itemize}

\section{Case Study}

\subsection{Calculator}

\begin{lstlisting}
send calc.add(3,4)
\end{lstlisting}

\subsection{AI Pipeline}

\[
User \rightarrow Planner \rightarrow LLM \rightarrow Memory
\]

\section{Discussion}

Key insight:

\[
\text{Value typing} \neq \text{Communication correctness}
\]

\section{Related Work}

Actor systems, session types, and AIOS frameworks.

\section{Conclusion}

We proposed AICE and demonstrated that type safety alone is insufficient for asynchronous actor systems.

Future work includes protocol inference, formal verification, and deployment on embedded systems.

\bibliographystyle{plain}
\bibliography{references}

\end{document}
